XElements extends GenFunctions to vendor-specific elements. The methods generate
basic elements such as gates and flip-flops.

The IBUF Element provides a simple example -
\begin{lstlisting}[gobble=0, language=java]
    public Element IBUF (Net o, Net i) {
        Element c = new Element("IBUF");
        c.addInput("I", i);
        c.addOutput("O", o);
        return(c);
    }
\end{lstlisting}
The new Element is created by a constructor, the element name
string being passed as an argument to the constructor. The input
and output Nets are added paired with the port name strings. The
Element is returned.

Where an Element may be subject to later optimisation a gate type
is added, e.g. the XORCY Element -
\begin{lstlisting}[gobble=0, language=java]
    public Net XORCY (Net li, Net ci) {
        Element c = new Element("XORCY", Gtype.XORCY);
        Net     o = new Net();
        c.addInput("LI", li);
        c.addInput("CI", ci);
        c.addOutput("O", o);
        return(o);
    }
\end{lstlisting}
In this case the output Net is returned.

The following basic Elements are constructed directly -

\blist
\item   LUT - a 1 to n input look-up table
\item   INV - inverter
\item   AND - 1 to n input AND gate
\item   OR - 1 to n input OR gate
\item   XOR - 1 to n input XOR gate
\item   XNOR - 1 to n input XNOR gate
\item   SRL16E - 16 bit CLB shift register
\item   SRL32E - 32 bit CLB shift register
\item   FDRS - reset/set flip-flop
\item   FDCP - clear/preset flip-flop
\item   MULT\_AND - multiplier AND
\item   MUXCY - carry multiplexer
\item   MUXF5 - CLB multiplexer
\item   MUXF6 - CLB multiplexer
\item   XORCY - carry XOR
\item   BUFG - global clock buffer
\item   IPAD - input pad
\item   OPAD - output pad
\item   IOPAD - bi-directional pad
\item   BUFE - internal 3-state buffer, +ve enable
\item   BUFT - internal 3-state buffer, -ve enable
\item   IBUF - input buffer
\item   IBUFDS - differential input buffer
\item   IBUFG - input global clock buffer
\item   IBUFGDS - differential input global clock buffer
\item   OBUF - output buffer
\item   OBUFDS - differential output buffer
\item   PULLDOWN - internal longline pulldown
\item   PULLUP - internal longline pullup
\item   RAM\_1 - single port CLB RAM
\item   ROM\_1 - single port CLB ROM
\item   RAM\_2 - dual port CLB RAM
\blend

Other Elements available in some Xilinx FPGA families but not others
have methods which must be overridden in subclasses -

\blist
\item   MULT18X18 - 18 x 18 bit multiplier
\item   MULT18X18S - 18 x 18 bit multiplier with output register
\item   MULT25X18 - 25 x 18 bit multiplier
\item   MULT25X18S - 25 x 18 bit multiplier with output register
\item   MULREG - multiplier used as register only
\item   ramallocate - block RAM
\item   fifoallocate - block RAM as FIFO
\item   FIFO - FIFO using dual port block RAM
\item   Mult- multiplier
\item   dsp48\_alu - an ALU using a DSP block
\blend
These are overridden in subclasses XC2V, XC2VP etc. where appropriate.

This subdivision is somewhat arbitrary; some Elements in the first
group are not available in some FPGAs or CPLDs and conditional code
ensures that these are not called is such cases. For consistency some
of these should be moved to the second group and the methods moved
into subclasses.

